\begin{tabularx}{\textwidth}{L{31mm}X X}
\toprule
\textbf{Modell} & \textbf{Lényege} & \textbf{Jellemző probléma} \\
\midrule
Egyprogramos modell & A RAM egyik részén az operációs rendszer, a másikon egyetlen felhasználói program található. & Egyszerű, de nincs valódi multiprogramozás, és pazarló lehet. \\
Fix partíciók & A memória előre rögzített méretű blokkokra van osztva, egy folyamat egy partíciót kap. & Belső töredezettség és folyamatméret-korlát jelentkezhet. \\
Változó partíciók & A partíciók mérete a folyamatok igényei szerint alakul ki. & Csökkenti a belső töredezettséget, de külső töredezettséget okozhat. \\
Lapozás & A logikai és fizikai memória azonos méretű lapokra, illetve keretekre bomlik. & Külső töredezettség helyett laptábla-kezelési költség és laphiba jelenik meg. \\
Szegmentáció & A program logikai egységekre, különböző méretű szegmensekre bomlik. & Jól illeszkedik a program szerkezetéhez, de külső töredezettséget okozhat. \\
\bottomrule
\end{tabularx}
